\section{Residual-treatment and conditional-information support}\label{app:support}

Support depends on the statement being estimated. A continuous horse race consumes residual score variation after controls; a tail comparison consumes weighted group membership; a ranking statement consumes pairwise order. The diagnostics below are therefore not generic data-coverage statistics.

For the primary beta tail contrast, estimator information is concentrated more sharply than the nominal count suggests. The frozen information audit reports an effective occupation count of approximately 43.3 for the fitted Q5-versus-Q1 contrast, compared with 468 nominal occupations. The pre-outcome residual-treatment audit reports an effective count of 53.263. These objects have different weights and should not be equated.

The computerization comparison uses Webb's software-exposure measure as the frozen primary comparison technology. Archived O*NET computer-use measures, routine-task intensity, and Frey--Osborne automation risk are sensitivities. Remotability is a separate proxy and not a measure of computerization. The comparison-technology results show that coefficient magnitude depends on the conditioning estimand; they do not select a uniquely correct exposure measure.

\begin{table}[htbp]
\centering
\caption{Effective residual support across comparison technologies}
\label{tab:app-residualsupport}
\scriptsize
\setlength{\tabcolsep}{3pt}
\begin{threeparttable}
\begin{tabularx}{\textwidth}{Yrrrrr}
\toprule
Architecture & Webb & O*NET importance & O*NET level & RTI & Frey--Osborne\\
\midrule
AIOE administrative & 72.2 & 73.0 & 77.3 & 62.8 & 63.8\\
AIOE ability-direct & 74.6 & 77.8 & 82.2 & 68.2 & 59.3\\
AIOE source-weighted & 72.1 & 65.4 & 70.3 & 62.1 & 62.8\\
Alpha & 17.4 & 31.1 & 31.7 & 11.9 & 26.3\\
Beta & 53.3 & 63.2 & 59.5 & 51.7 & 62.9\\
Broad & 84.5 & 44.4 & 72.2 & 77.5 & 74.0\\
\bottomrule
\end{tabularx}
\begin{tablenotes}[flushleft]\footnotesize
\item Notes: Inverse-Herfindahl counts of employment-weighted residual-exposure variance after projecting each exposure on the named comparison technology. Frozen pre-outcome treatment-side analysis; no labor outcome enters.
\end{tablenotes}
\end{threeparttable}
\end{table}

Top-five residual-variance shares for the same columns are 18.4, 17.9, 16.6, 20.5, and 20.0 percent for administrative AIOE; 15.9, 16.5, 15.0, 17.8, and 20.6 for ability-direct AIOE; 18.1, 19.5, 18.3, 20.6, and 20.2 for source-weighted AIOE; 41.6, 32.8, 33.3, 46.6, and 36.0 for alpha; 22.2, 20.1, 20.6, 22.8, and 19.8 for beta; and 15.8, 25.2, 18.0, 17.3, and 17.1 for broad exposure.

\subsection{Support after broad-family conditioning}

The repaired baseline uses 468 occupations in 22 SOC2 groups. Every group contains at least two beta quintiles, but only four contain both Q1 and Q5. Thus a Q5--Q1 coefficient can still be estimated after SOC2 conditioning through a connected set of intermediate-quintile comparisons, even though most families do not supply a direct tail-to-tail contrast.

Conditional target information falls from 26.44 million in the baseline to 7.95 million after adding SOC2-by-post terms and 7.93 million after absorbing SOC2-by-month paths. The corresponding weighted residual SD falls from .158 to .087. Effective-information counts rise from 43.36 to about 55.8 because the remaining information is less concentrated, but this should not be mistaken for more total identifying information: only 30 percent of the baseline information remains. The estimated coefficient consequently has an SE around .0705 rather than .0450.

The original six-score stable tails supply a different support diagnostic. Forty-six occupations are always Q1 and 18 always Q5, together representing 9.74 percent of employment. The stable-tail estimate has effective information 15.02 and a top-five share of 46.7 percent. Common occupational coverage therefore does not imply common classification or diffuse estimator support.
